% GENERATED FILE. Edit canonical lessons, then run npm run generate:books.
% canonical-source-hash: fnv1a64:6dabf95ea4eb2c18
% canonical-lessons: TE-C158-tondarapadu, TE-C158-tomu, TE-C158-marcu, TE-C158-vippu, TE-C158-kolucu

\chapter{To hurry, To brush, To change, To take off, To measure}
\label{ch:te-to-hurry-to-brush-to-change-to-take-off-to-measure}
\hlchaptermodality{158}

\begin{chapteropening}
\textbf{By the end of this chapter:} \emph{I can say to hurry, to brush, to change, to take off and to measure.}

Complete the last lesson of chapter 158: I can say to hurry, to brush, to change, to take off and to measure.
\end{chapteropening}

\section[\texorpdfstring{\te{తొందరపడు} (tondarapaḍu)}{తొందరపడు (tondarapaḍu)}]{\te{తొందరపడు} (tondarapaḍu) — to hurry}

\label{lesson:TE-C158-tondarapadu}



\begin{quote}
Before the new one: say the Telugu for to mix, then the Telugu for to grind.
\end{quote}

\subsection*{You'll want to know: \te{తొందరపడు}}
\textbf{\te{తొందరపడు}} — \emph{tondarapaḍu} — \textquotedblleft{}to hurry\textquotedblright{}.

\begin{grammarlens}[title={What you've built}]
One more word for this chapter.
\end{grammarlens}

\subsection*{Your turn}
\begin{itemize}
\raggedright
  \item \emph{Say it:} \emph{tondarapaḍu}
  \item \emph{Say it:} \emph{tondarapaḍu}, once more
  \item \emph{Say it:} say \emph{tondarapaḍu}
  \item \emph{Recall:} say the Telugu for to mix, then the Telugu for to grind, then say \emph{tondarapaḍu} again
\end{itemize}

\begin{tcolorbox}[breakable,colback=teal!4,colframe=teal!35!black,title={Before you move on}]
What does \te{తొందరపడు} mean? (\textquotedblleft{}To hurry\textquotedblright{}.) Say it once more.
\end{tcolorbox}
\section[\texorpdfstring{\te{తోము} (tōmu)}{తోము (tōmu)}]{\te{తోము} (tōmu) — to brush (teeth), to scour (pots)}

\label{lesson:TE-C158-tomu}



\begin{quote}
Before the new one: say the Telugu for to grind, then the Telugu for to hurry.
\end{quote}

\subsection*{You'll want to know: \te{తోము}}
\textbf{\te{తోము}} — \emph{tōmu} — \textquotedblleft{}to brush (teeth), to scour (pots)\textquotedblright{}.

\begin{grammarlens}[title={What you've built}]
One more word for this chapter.
\end{grammarlens}

\subsection*{Your turn}
\begin{itemize}
\raggedright
  \item \emph{Say it:} \emph{tōmu}
  \item \emph{Say it:} \emph{tōmu}, once more
  \item \emph{Say it:} say \emph{tōmu}
  \item \emph{Recall:} say the Telugu for to grind, then the Telugu for to hurry, then say \emph{tōmu} again
\end{itemize}

\begin{tcolorbox}[breakable,colback=teal!4,colframe=teal!35!black,title={Before you move on}]
What does \te{తోము} mean? (\textquotedblleft{}To brush (teeth), to scour (pots)\textquotedblright{}.) Say it once more.
\end{tcolorbox}
\section[\texorpdfstring{\te{మార్చు} (mārcu)}{మార్చు (mārcu)}]{\te{మార్చు} (mārcu) — to change}

\label{lesson:TE-C158-marcu}



\begin{quote}
Before the new one: say the Telugu for to hurry, then the Telugu for to brush.
\end{quote}

\subsection*{You'll want to know: \te{మార్చు}}
\textbf{\te{మార్చు}} — \emph{mārcu} — \textquotedblleft{}to change\textquotedblright{}.

\begin{grammarlens}[title={What you've built}]
One more word for this chapter.
\end{grammarlens}

\subsection*{Your turn}
\begin{itemize}
\raggedright
  \item \emph{Say it:} \emph{mārcu}
  \item \emph{Say it:} \emph{mārcu}, once more
  \item \emph{Say it:} say \emph{mārcu}
  \item \emph{Recall:} say the Telugu for to hurry, then the Telugu for to brush, then say \emph{mārcu} again
\end{itemize}

\begin{tcolorbox}[breakable,colback=teal!4,colframe=teal!35!black,title={Before you move on}]
What does \te{మార్చు} mean? (\textquotedblleft{}To change\textquotedblright{}.) Say it once more.
\end{tcolorbox}
\section[\texorpdfstring{\te{విప్పు} (vippu)}{విప్పు (vippu)}]{\te{విప్పు} (vippu) — to take off, to untie}

\label{lesson:TE-C158-vippu}



\begin{quote}
Before the new one: say the Telugu for to brush, then the Telugu for to change.
\end{quote}

\subsection*{You'll want to know: \te{విప్పు}}
\textbf{\te{విప్పు}} — \emph{vippu} — \textquotedblleft{}to take off, to untie\textquotedblright{}.

\begin{grammarlens}[title={What you've built}]
One more word for this chapter.
\end{grammarlens}

\subsection*{Your turn}
\begin{itemize}
\raggedright
  \item \emph{Say it:} \emph{vippu}
  \item \emph{Say it:} \emph{vippu}, once more
  \item \emph{Say it:} say \emph{vippu}
  \item \emph{Recall:} say the Telugu for to brush, then the Telugu for to change, then say \emph{vippu} again
\end{itemize}

\begin{tcolorbox}[breakable,colback=teal!4,colframe=teal!35!black,title={Before you move on}]
What does \te{విప్పు} mean? (\textquotedblleft{}To take off, to untie\textquotedblright{}.) Say it once more.
\end{tcolorbox}
\section[\texorpdfstring{\te{కొలుచు} (kolucu)}{కొలుచు (kolucu)}]{\te{కొలుచు} (kolucu) — to measure}

\label{lesson:TE-C158-kolucu}



\begin{quote}
Before the new one: say the Telugu for to change, then the Telugu for to take off.
\end{quote}

\subsection*{You'll want to know: \te{కొలుచు}}
\textbf{\te{కొలుచు}} — \emph{kolucu} — \textquotedblleft{}to measure\textquotedblright{}.

\begin{grammarlens}[title={What you've built}]
One more word for this chapter.
\end{grammarlens}

\subsection*{Your turn}
\begin{itemize}
\raggedright
  \item \emph{Say it:} \emph{kolucu}
  \item \emph{Say it:} \emph{kolucu}, once more
  \item \emph{Say it:} say \emph{kolucu}
  \item \emph{Recall:} say the Telugu for to change, then the Telugu for to take off, then say \emph{kolucu} again
\end{itemize}

\begin{tcolorbox}[breakable,colback=teal!4,colframe=teal!35!black,title={Before you move on}]
What does \te{కొలుచు} mean? (\textquotedblleft{}To measure\textquotedblright{}.) Say it once more.
\end{tcolorbox}
